\chapter{参考书对照与公式核查}\label{app:references}
\section{本稿与两本主要参考书的对应}
下列对照依据本地扫描 PDF 的目录与所检查的正文。印刷页码指原书页面中的页码；PDF 页序从 1 计数，包括前置页。两者切勿混用。
本附录专用于第二部分“广义函数专题扩展”；吉田主线的版本对照见附录 \ref{app:yosida-map}。

\begin{center}\small
\begin{tabular}{p{0.25\linewidth}p{0.43\linewidth}p{0.22\linewidth}}
\toprule
本稿主题 & 原书位置 & 页码线索\\\midrule
动机、分布定义 & 施瓦尔茨，第 I 章 §§1--3 & 印刷 1--16；章首 PDF 24\\
测试空间拓扑 & 施瓦尔茨，第 III 章 §§1--3 & 印刷 43--53；章首 PDF 66\\
求导与局部结构 & 施瓦尔茨，第 II 章；第 III 章 §§4--10 & 第 II 章 PDF 40；点支集印刷 70 / PDF 93\\
张量积与光滑乘法 & 施瓦尔茨，第 IV--V 章 & 章首 PDF 96、106\\
卷积与基本解 & 施瓦尔茨，第 VI 章 §§2、10 & 卷积印刷 111 / PDF 134\\
温和分布与 Fourier & 施瓦尔茨，第 VII 章 §§2--6 & 印刷 169、172 / PDF 192、195\\
凸权重与 $W$ 型空间 & 盖尔芳特、希烙夫，第 I 章 §§1--3 & 印刷 1--22 / PDF 5--26\\
柯西问题与适定性 & 盖尔芳特、希烙夫，第 II--III 章 & 章首 PDF 30、107\\
广义固有函数 & 盖尔芳特、希烙夫，第 IV 章 §§4--6 & 印刷 177--198 / PDF 181--202\\
\bottomrule
\end{tabular}
\end{center}
盖尔芳特、希烙夫扫描目录的第一章标题显示为 $M$ 型空间，正文标题使用 $W$ 型空间，故本稿优先遵从正文符号。该卷本身不是广义函数初步教程，它多处调用前两卷结果。本稿将其作为应用和拓展的主要参考，基础内容以施瓦尔茨的系统叙述为主。

\section{运算条件速查}
\begin{center}\small
\begin{tabular}{p{0.23\linewidth}p{0.43\linewidth}p{0.24\linewidth}}
\toprule
运算 & 本稿使用的充分条件 & 所得对象\\\midrule
$\partial^\alpha T$ & $T\in\Dist(\Omega)$ & 分布\\
$aT$ & $a\in C^\infty(\Omega)$ & 分布\\
$T*\rho$ & $T\in\Dist(\R^n),\rho\in\Test$ & 光滑函数\\
$S*T$ & 至少一个因子紧支集 & 分布\\
$S*T$ & 加法在支集乘积上适当 & 分布\\
$\widehat T$ & $T\in\Temp(\R^n)$ & 温和分布\\
$S*T$ 的 Fourier 公式 & $S\in\CompactDist,T\in\Temp$ & 温和分布\\
$F^*T$ & $F$ 是光滑微分同胚 & 分布\\
$G_t*T$ & $t>0,T\in\Temp$ & 光滑函数与温和分布\\
\bottomrule
\end{tabular}
\end{center}
这些都是明确可用的充分条件，不宣称列尽所有可能的合法情形。

\section{检查证明时的五个问题}
\begin{enumerate}
\item 测试空间是什么？要求共同紧支集，还是全空间带权衰减？
\item 配对是否线性？是否混入了 Hilbert 内积的复共轭？
\item 对偶运算由哪个测试函数映射的转置得到？该映射是否连续？
\item 移动导数、极限或积分时，支集和半范数是否有统一控制？
\item 唯一性和稳定性在哪个解空间与数据拓扑中成立？
\end{enumerate}

\Needspace{14\baselineskip}
\section{与已完成主线的衔接}
测试空间与紧支集分布衔接第一章 §§8、12--14；分布极限衔接第二章 §3；局部结构衔接第三章 §11。Fourier、卷积及 Sobolev 估计与第六章 §§1--9 互补；本稿第六章 §4 已有 Paley--Wiener 内容，不再把它整体列为未写专题。紧嵌入、Dirichlet 弱解和核空间分别见第十章 §§3、6 及其附录。一般自伴谱测度在第十一章，半群及发展方程在第九、十四章。

本轮新增的全空间嵌入、带质量项椭圆模型和加权空间稠密性，分别见命题 \ref{prop:topic-sobolev-embedding}、\ref{prop:topic-bessel-inverse}、\ref{prop:weighted-density}；边界迹与反向密度演化见书后补充说明。本稿尚未证明一般 Schwartz 核定理、完整 $W$ 型空间 Fourier 同构、一般装配空间谱定理以及非光滑区域的完整迹与正则性定理。查阅这些主题时应保留各节注明的外部前提，不能把有限模型的证明提升为一般结论。
